% !TeX root = main.tex
  
\section{Naive Bayes}

In statistics, naive Bayes classifiers are a family of simple "probabilistic classifiers" based on applying Bayes' theorem with strong (naive) independence assumptions between the features. They are among the simplest Bayesian network models, but coupled with kernel density estimation, they can achieve high accuracy levels.\footnote{\url{https://en.wikipedia.org/wiki/Naive_Bayes_classifier}} Basically assume the features are independent, for example if classifying fruits apples vs pears, we assume the size of the pear or apple is completely independent of the color, if we think of color and size as the features used in classification.\autocite{hastie2009,bishop2006}

\subsection{Independency Assumptions}

 \begin{itemize}
    \item still widely (and successfully) used 
    \item often outperforming much more advanced classifiers 
    \item appropriate in the presence of high dimensional features (curse of dimensionality). it is used to overcome the curse of dimensionality, for eg. an image will have all the pixels as the dimensions which might be 1024x1024 and that vector is too large !! 
    \item also called \textbf{''Idiot's Bayes''}
   \end{itemize}
   
   \textbf{Naive Bayes and Statistical Independency}
  
  For the class-conditional pdf $p(\vect x\mid y)$ we can do the following factorization:
%
  \begin{eqnarray*}
    p(\vect x\mid y) &=& p(x_1, x_2, \dots, x_d \mid y) \\
    &=& p(x_1\mid y)p(x_2, x_3,\dots, x_d \mid y, x_1) \\
    &=& p(x_1\mid y)p(x_2\mid y,x_1)p(x_3,x_4,\dots, x_d \mid y, x_1, x_2) \\
    &=& p(x_1\mid y)\prod_{i=2}^d p(x_i\mid y, x_1,\dots,x_{i-1})
  \end{eqnarray*} 
   \begin{itemize}
    \item The Naive Bayes classifier makes a strong \textbf{naive independence assumption}.
    \item All $d$ components $x_i$ of the feature vector $\vect x$ are assumed to be \textbf{conditionally} independent given the class $y$.
    \item This implies:
      \begin{eqnarray*}
        p(\vect x\mid y) &=& \prod_{i=1}^d p(x_i\mid y)
      \end{eqnarray*}
  \end{itemize}

  \textbf{Decision rules:}
  \begin{itemize}
    \item Bayes-optimal prediction: $\hat y^* = \argmax_y p(y\mid\vect x)$.
    \item Naive Bayes prediction (using the conditional-independence assumption):
      \begin{eqnarray*}
        \hat y &=& \argmax\limits_{y} p(y)\prod_{i=1}^d p(x_i \mid y)
      \end{eqnarray*}
      In general, $\hat y$ is only an approximation of $\hat y^*$.
  \end{itemize}

\subsection{An Example: Gaussians}
  Assume the $100$--dimensional feature vector $\vect x\in \real^{100}$ belonging to class $y$ is normally distributed and all components are {\em mutually dependent}:
 
    \begin{eqnarray*}
      \vect\mu_y &\in& \real^{100} \\
      \mat\Sigma &=&  \mat\Sigma^T \in \real^{100\times 100} 
    \end{eqnarray*}
    The total number of parameters to be estimated for each class is  $$100+100\cdot (100+1)/2= 5150.$$
 Assume the $100$--dimensional feature vector $\vect x\in \real^{100}$ belonging to class $y$ is normally distributed and all components are \textbf{mutually independent}. 
    \begin{displaymath}
      p(\vect x\mid y) = \prod_{i=1}^{100} p(x_i \mid y) \quad = \quad \prod_{i=1}^{100} \mathcal N(x_i;\mu_{y,i}, \sigma^2_{y,i}).
    \end{displaymath}
    For each component $i\in\{1,\dots,100\}$ we have to estimate mean $\mu_{y,i}\in \real$ and variance $\sigma_{y,i}^2\in \real$.
    The total number of parameters to be estimated for each class is  $$100+100= 200.$$\\
    The slides show examples of quadratic decision boundaries considering the statistical dependency and with independence assumption of the variables $x1 and x2$. 
    \textbf{Exam Notes:} Important to understand the difference clearly in these cases. \\
   Let us consider the \textbf{logit transform} for two classes $y\in\{0,1\}$:
   \begin{eqnarray*}
    \log \frac{p(y=0\mid\vect x)}{p(y=1\mid\vect x)}
      &=& \log\frac{p(y=0) p(\vect x \mid y=0)}{p(y=1) p(\vect x \mid y=1)}\\
      &=& \log\frac{p(y=0)}{p(y=1)} + \log\frac{p(\vect x \mid y=0)}{p(\vect x \mid y=1)}\\
      &=& \log\frac{p(y=0)}{p(y=1)} + \log\frac{\prod_{i=1}^d p(x_i \mid y=0)}{\prod_{i=1}^d p(x_i \mid y=1)} \\
      &=& \underbrace{f_0 + \sum_{i=1}^d f_i(x_i)}_{\mbox {\textbf{generalized additive model}}}
   \end{eqnarray*}
   with the constant $f_0=\log\frac{p(y=0)}{p(y=1)}$ and one function per feature, $f_i(x_i)=\log\frac{p(x_i\mid y=0)}{p(x_i\mid y=1)}$. This is the log-odds $F(\vect x)$ of Chapter~03. (Different symbols on purpose: in Chapter~03, $\alpha_0$ also contains the mean and covariance terms and $\alpha_i$ are constant coefficients.)
   Understand the difference between Naive Bayes and Bayes. There are multiple techniques to beat the curse of dimensionality, for example:
   \begin{itemize}
     \item Reduction of the parameter space
       \begin{itemize}
         \item Introduction of independency assumptions
           (from complete dependency to mutual independency)
         \item Parameter tying
       \end{itemize}
    \item Reduction of the dimension of the feature vectors       
  \end{itemize}
   First order dependency
  \begin{eqnarray*}
    p(\vect x\mid y) &=& p(x_1, x_2, \dots, x_d \mid y) \\
                 &=& p(x_1\mid y)p(x_2, x_3,\dots, x_d \mid y, x_1) \\
                 &=& p(x_1\mid y)p(x_2\mid y,x_1)p(x_3,x_4,\dots, x_d \mid y, x_1, x_2) \\
                 &=& p(x_1\mid y)\prod_{i=2}^d p(x_i\mid y, x_1,\dots,x_{i-1}) \\
                 &\approx& p(x_1\mid y)\prod_{i=2}^d p(x_i\mid y, x_{i-1})
  \end{eqnarray*}
  where the last step is the first-order dependency assumption $p(x_i\mid y,x_1,\dots,x_{i-1})\approx p(x_i\mid y,x_{i-1})$.
  Example:
    \textbf{First order dependency} in a Gaussian random vector corresponds to a tridiagonal structure in the \textbf{inverse covariance (precision) matrix} $\mat Q = \mat\Sigma^{-1}$ (in general $\mat\Sigma$ itself is dense):

    \begin{eqnarray*}
      \mat Q &=& \left(
        \begin{array}{ccccccc}
          q_{1,1} & q_{2,1} & 0            & 0            & \cdots & 0              & 0 \\
          q_{2,1} & q_{2,2} & q_{3,2} & 0            & \cdots & 0              & 0 \\
          0            & q_{3,2} & q_{3,3} & q_{4,3} & \cdots & 0              & 0 \\
          0            & 0            & q_{4,3} & q_{4,4} & \cdots & 0              & 0 \\
          \vdots       & \vdots       & \vdots       & \vdots       & \ddots & \vdots         & q_{d,d-1} \\
          0            & 0            & 0            & 0            & \cdots & q_{d,d-1} & q_{d,d} \\
        \end{array}
      \right)
    \end{eqnarray*}
    Example: first-order dependency with \textbf{tied diagonal elements}, i.\,e.\ $q_{i,i}=q$:
      \begin{eqnarray*}
     \mat Q &=& \left(
       \begin{array}{ccccccc}
         q       & q_{2,1} & 0            & 0            & \cdots & 0              & 0 \\
         q_{2,1} & q       & q_{3,2} & 0            & \cdots & 0              & 0 \\
         0            & q_{3,2} & q       & q_{4,3} & \cdots & 0              & 0 \\
         0            & 0            & q_{4,3} & q       & \cdots & 0              & 0 \\
         \vdots       & \vdots       & \vdots       & \vdots       & \ddots & \vdots         & q_{d,d-1} \\
         0            & 0            & 0            & 0            & \cdots & q_{d,d-1} & q \\
       \end{array}
     \right)
    \end{eqnarray*}
    
    \subsection{Lessons Learned}
\begin{itemize}
    \item Na{\"i}ve Bayes is rather successful. 
    \item Na{\"i}ve Bayes does not require a huge set of training data. 
    \item Statistical dependency vs. dimension of the search space. 
    \item Na{\"i}ve Bayes: give it a try!
  \end{itemize}

\subsection{Comprehensive Questions}
 \begin{itemize}
    \item What is the assumption of Naive Bayes? 
    \item How does the assumption affect the class-conditional pdf $p(\vect x\mid y)$? 
    \item What is the structure of the covariance matrix of normal-distributed classes in Naive Bayes? 
    \item How can Naive Bayes be extended to first-order statistical dependencies?
  \end{itemize}
  \newpage

\section{Discriminative Analysis}

 \textbf{Discriminant analysis} methods are \textbf{generative modeling} methods: they model the prior $p(y)$ and the class-conditional density $p(\vect x\mid y)$ and obtain the posterior through Bayes' rule.\autocite{hastie2009,bishop2006}
  
  \begin{displaymath}
    p(y\mid\vect x) = \frac{p(y)\cdot p(\vect x\mid y)}{\sum_{y'} p(y')\cdot p(\vect x\mid y')}
  \end{displaymath}
  LDA is a generative method despite its name, since it models the joint probability of the data and class.\footnote{https://stats.stackexchange.com/questions/411824/linear-discriminant-analysis-generative-or-discriminative}
  
  \subsection{Gaussian Classifier}
   We call the Bayesian classifier \textbf{Gaussian}, if the class conditional density $p(\vect x|y)$ is Gaussian,(for eg. the availability of apples is gaussian in nature throughout the year) i.\,e.

  \begin{eqnarray*}
     p(\vect x | y) &=& \mathcal {N} (\vect x ;\vect {\mu}_y, \mat\Sigma_y) \\
                   &=& \frac{1}{\sqrt{\det 2\pi  \mat\Sigma_y}} 
                       e^{-\frac{1}{2}(\vect x -\vect{\mu}_y)^T \mat\Sigma_y^{-1}(\vect x - \vect{\mu}_y)}
  \end{eqnarray*}
  %
  where
%
  \begin{center}
    \begin{minipage}{0.6\textwidth}
      \begin{itemize}
        \item[$\vect x \in \real^d$:] $d$-dimensional feature vector\\
        \item[$\vect {\mu}_y \in \real^d $:] mean vector of class y\\
        \item[$ \mat\Sigma_y \in \real^{d\times d}$:] positive definite covariance matrix.
      \end{itemize}
    \end{minipage}
  \end{center}
  
  \textbf{Facts about Gaussian classifiers:}

  \begin{itemize}
    \item In general the decision boundary is \textbf{quadratic} in the components $x_j$ 
      of the feature vector $\vect x$. 
    \item If all classes share the same covariance, the decision boundary is \textbf{linear} in the components $x_j$ of the feature vector $\vect x$. 
    \item If all covariance matrices are diagonal matrices, then we get a \textbf{Naive Bayes} classifier.
    \item If all classes share the same (pooled) covariance matrix $\mat\Sigma$ and priors are identical, the optimal (Bayes) decision minimizes the \textbf{Mahalanobis distance}
      \begin{eqnarray*}
        \hat y^* &=& \argmin_{y} \frac{1}{2} (\vect x-\vect\mu_y)^T\mat\Sigma^{-1}(\vect x-\vect\mu_y)
      \end{eqnarray*}
    \item If all covariance matrices are the identity matrix, we get the \textbf{nearest class mean} classifier based on the $L_2$-norm:
      \begin{eqnarray*}
        \hat y^* &=& \argmin_{y} \frac{1}{2} (\vect x-\vect\mu_y)^T(\vect x-\vect\mu_y)
      \end{eqnarray*}
      The prototype vectors are the mean vectors.
        
  \textbf{From linear to quadratic decision boundaries:} 
  A compromise between linear and quadratic decision boundaries can be achieved by using \textbf{regularized covariance matrices}:
%  
  \begin{eqnarray*}
    \mat\Sigma_y (\alpha)= \alpha \mat\Sigma_y + (1-\alpha)\mat\Sigma
  \end{eqnarray*}  
  where $\alpha\in [0,1]$ and $\mat\Sigma$ denotes the shared (pooled) covariance.  Obviously we have the extremes:
  \begin{itemize}
    \item Linear decision boundary: $\quad ~~\alpha = 0$
    \item Quadratic decision boundary: $\alpha =1$
  \end{itemize}
  \end{itemize}
  
 \subsection{Feature Transform}
  Can we find a feature transform
%  TODO fix this slide
  \begin{eqnarray*}
    \phi: \real^d &\rightarrow &\real^d
   \end{eqnarray*}
  to generate features $\phi(\vect x)$ that share the same covariance matrix?  This is how we do it -\\
   The symmetric positive definite covariance matrix $\mat\Sigma\in \real^{d\times d}$ can be decomposed using SVD (here identical to the eigen decomposition).\autocite{trefethen1997,golub1996}
%  
  \begin{eqnarray*}
    \mat\Sigma&=&  \mat{U}\mat{D}\mat{U}^T = (\mat{U}\mat{D}^{\frac{1}{2}})(\mat{U}\mat{D}^{\frac{1}{2}})^T  =      (\mat{U}\mat{D}^{\frac{1}{2}})\cdot \mat{I} \cdot(\mat{U}\mat{D}^{\frac{1}{2}})^T
  \end{eqnarray*}
  where $\mat{I}\in \real^{d\times d}$ is the identity matrix.
 
  \begin{itemize}
    \item \textbf{Determinant:}  
      \begin{displaymath}
        \det \mat\Sigma = \prod_{i=1}^d \lambda_i,
      \end{displaymath}
      where $\mat D=\mathrm{diag}(\lambda_1,\dots,\lambda_d)$ and $\lambda_i>0$ are the \textbf{eigenvalues} (= singular values) of $\mat\Sigma$.
    \item \textbf{Inverse:}
      \begin{displaymath}
        \mat\Sigma^{-1} = \mat{U}\mat{D}^{-1}\mat{U}^T
                          = (\mat{U}\mat{D}^{-\frac{1}{2}})\cdot \mat{I} \cdot(\mat{U}\mat{D}^{-\frac{1}{2}})^T
      \end{displaymath}
  \end{itemize}
     Now we incorporate this:
   
  \begin{eqnarray*}
    \mathcal {N} (\vect x ;\vect {\mu}, \mat\Sigma) 
      &=& \frac{1}{\sqrt{\det 2\pi  \mat\Sigma}} 
          e^{-\frac{1}{2}(\vect x -\vect{\mu})^T \, \mat\Sigma^{-1} \, (\vect x - \vect{\mu})} \\
      &=& \frac{1}{\sqrt{\det 2\pi  \mat\Sigma}} 
          e^{-\frac{1}{2}(\vect x -\vect{\mu})^T \, (\mat{U}\mat{D}^{-\frac{1}{2}})\cdot \mat{I} \cdot(\mat{U}\mat{D}^{-\frac{1}{2}})^T \, (\vect x - \vect{\mu})} \\
      &=& \frac{1}{\sqrt{\det 2\pi  \mat\Sigma}} 
          e^{-\frac{1}{2} \big( (\mat{D}^{-\frac{1}{2}}\mat{U}^T)\vect x - (\mat{D}^{-\frac{1}{2}}\mat{U}^T)\vect\mu \big)^T \,\mat{I} ~
          \big( (\mat{D}^{-\frac{1}{2}}\mat{U}^T)\vect x- (\mat{D}^{-\frac{1}{2}}\mat{U}^T)\vect\mu \big) }
  \end{eqnarray*}
  The classwise transform $\phi_y$ is even a linear function:
%
  \begin{displaymath}
    \vect x' \quad = \quad \phi_y(\vect x) \quad = \quad \mat{D}_y^{-\frac{1}{2}}\mat{U}^T_y\vect x
  \end{displaymath}
  It is straight forward  to show that $\vect x'$ is normally distributed
  \begin{displaymath}
     p(\vect x' | y) \quad = \quad \mathcal {N} (\vect x' ; \vect {\mu}'_y, \mat\Sigma'_y) 
                    \quad = \quad \mathcal {N} (\vect x' ; \mat{D}_y^{-\frac{1}{2}}\mat{U}^T_y\vect {\mu}_y, \mat{I}) 
  \end{displaymath}
   \textbf{Conclusions:}
  
  \begin{itemize}
    \item All classes $y$ share the same covariance matrix that is the identity matrix.
    \item With class-specific transforms $\phi_y$, the boundary in the original feature space is still in general \emph{quadratic}; it becomes linear (a \emph{hyperplane}) only in the transformed space for each fixed $y$.
    \item \textbf{Huge disadvantage:} feature transform depends on the class index $y$ (class-specific transform $\phi_y$). This is undesirable: for large data sets we would need to estimate and apply a separate transform for each class.
    \item If we have a classified training set, we can compute a transform for each class such that all covariance matrices are the identity matrix.
    \item Classification requires the application of different transforms.
  \end{itemize}
  
  \subsection{Linear Discriminant Analysis}
  
  A method to do the feature transform that is common for all samples. The idea is similar to before and the steps to do the actual feature transform is shared here. Note: need to understand linear algebra fundamentals on what is SVD, many resources available online. LDA is a generative method, despite its name.\footnote{\url{https://towardsdatascience.com/fishers-linear-discriminant-intuitively-explained-52a1ba79e1bb}} LDA is built on Fishers discriminant analysis which is explained later.\footnote{\url{https://www.geeksforgeeks.org/ml-linear-discriminant-analysis/}}\footnote{\url{https://en.wikipedia.org/wiki/Linear_discriminant_analysis}}
  
   \begin{algorithmic}
    \STATE \textbf{Input:} training data:  $S = \{ (\vect x_1, y_1), (\vect x_2, y_2), (\vect x_3, y_3), \dots, (\vect x_m, y_m) \}$\\
    \STATE 1. ML estimation of the \textbf{class means} and the \textbf{shared (pooled) covariance matrix}:
      \begin{displaymath}
        \hat{\vect\mu}_y = \frac{1}{m_y}\sum_{i:\,y_i=y} \vect x_i, \qquad
        \widehat{\mat\Sigma}= \frac{1}{m} \sum_{i=1}^m (\vect x_i - \hat{\vect \mu}_{y_i})(\vect x_i - \hat{\vect \mu}_{y_i})^T
      \end{displaymath}
      where $m_y = |\{i: y_i=y\}|$ is the number of training samples of class $y$.
     \STATE 2. Compute SVD of the estimated covariance matrix: $\widehat{\mat \Sigma}= \mat{U}\mat{D}\mat{U}^T$ 
     \STATE 3. Assign transform: 
       \begin{displaymath}
         \phi(\vect x)= \mat D^{-\frac{1}{2}} \mat{U}^T\vect x
       \end{displaymath}
     \STATE 4. Compute transformed mean vectors for all $y$
       \begin{displaymath}
          \vect{\mu}'_y = \phi(\hat{\vect\mu}_y)= \mat{D}^{-\frac{1}{2}} \mat{U}^T \hat{\vect{\mu}}_y
       \end{displaymath}
     \STATE \textbf{Output:} feature transform $\phi$, transformed mean vectors $\vect{\mu}'_y$
  \end{algorithmic}
   Decision rule using \textbf{sphered data $\phi(\vect x)$}: the predicted label $\hat y$ of a feature vector $\vect x$ is
  \begin{eqnarray*}
    \hat y &=&  \argmax_y p(y\mid \phi(\vect x))\\ 
        &=& \argmax_y 
              \left\{ 
                \log p(y) - \frac{1}{2} \big( \phi(\vect x) - \phi(\hat{\vect\mu}_y) \big)^T \big( \phi(\vect x)-\phi(\hat{\vect\mu}_y) \big)
              \right\} \\ 
        &=& \argmin_y \left\{ \frac{1}{2}\left\|\phi(\vect x)-\phi(\hat{\vect\mu}_y) \right\|_2^2 - \log p(y) \right\}
  \end{eqnarray*}
  where $\|.\|_2$ denotes the $L_2$ norm and the class means are replaced by their estimates $\hat{\vect\mu}_y$.\\
  (Plug-in classifier: $\hat y$ approximates the Bayes decision $\hat y^*$.)\\[0.3cm]
  
  \textbf{Conclusions:}\\
  \begin{itemize}
    \item If all classes share the \textbf{same prior},  the decision rule is the \textbf{nearest class mean} decision rule, where transformed mean vectors serve as prototypes.
    protoypes : refer to the mean of the clusters, could be scalar or vector depending on the samples x.
    \item The feature transform $\phi$ does not change the dimension of features.
  \end{itemize}
   \textbf{2 classes} $y\in\{0,1\}$: insights from geometrical analysis of sphered data. the most common case is to have only two classes, so throughout the document, we focus on 2 class cases. It is clear from the slides, that in the 2 class case, the angle can tell us if the sample belongs to one class or the other. so just 1D information can seperate the samples in a 2 class case. 
   
  \begin{itemize}
    \item Angle between $\phi(\vect x)$ and $(\phi(\hat{\vect \mu}_1)-\phi(\hat{\vect \mu}_0))$ can be used 
      for decision making.
    \item Decision rule (equal priors):
      {\small
      \begin{displaymath}
        \hat y = \left\{
                \begin{array}{cl}
                  0, & \quad \mbox{if} \quad 
                       \phi(\vect x)^T \big( \phi(\hat{\vect\mu}_1) - \phi(\hat{\vect\mu}_0) \big) < 
                       \frac{1}{2} \big( 
                         \phi(\hat{\vect{\mu}}_1)^T \phi(\hat{\vect{\mu}}_1) - 
                         \phi(\hat{\vect{\mu}}_0)^T \phi(\hat{\vect{\mu}}_0) 
                        \big) \\
%                       \big( \phi(\vect{\mu}_1) + \phi(\vect{\mu}_0) \big)^T
%                       \big( \phi(\vect{\mu}_1) - \phi(\vect{\mu}_0) \big) \\[.3cm]
                  1, & \quad \mbox{otherwise}.
                \end{array}
              \right.
      \end{displaymath}
      }
    \item Coordinate orthogonal to the 1-D subspace spanned by $(\phi(\hat{\vect \mu}_1)-\phi(\hat{\vect \mu}_0))$ does not affect relative distances.
  \end{itemize}
    \textbf{$K$ classes:} insights from geometrical analysis of sphered data. This is the multi-class case, better to understand as well. naturally derive from the previous case. in a $(K)$ class case, a $(K-1)$ dimension subspace can separate the samples between the $(K)$ classes. any other information is extra and not needed. 
  \begin{itemize}
    \item Class centroids span $(K-1)$-dimensional subspace. 
    \item Relative differences are not affected by coordinates in the $(d-K+1)$-dimensional subspace that is orthogonal to the  $(K-1)$-dimensional subspace spanned by class centroids.
  \end{itemize}
    \textbf{Objective:} \\
  Will we gain an advantage if we transform features by $$\phi: \real^d\rightarrow \real^k$$ in higher $(k>d)$ or lower dimensional $(k<d)$ spaces? This is something to think about and explored in detail in the next chapters. 
  
  \subsection{Lessons Learned}
\begin{itemize}
    \item Relationship between Bayesian classifier, Gaussian classifier, and nearest class mean classifier. 
    \item Mahalanobis distance 
    \item Linear Discriminant Analysis is a regularized nearest class mean classifier 
    \item Class centroids span $(K-1)$-dimensional subspace 
  \end{itemize}
  
  \subsection{Comprehensive Questions}
  \begin{itemize}
    \item What is a Gaussian classifier? 
    \item What is the idea behind the feature transform for the LDA? 
    \item Formulate the LDA for normally distributed classes. 
    \item What is the dimensionality of the LDA subspace for $K$ classes?
  \end{itemize}
  
  \subsection{Rank Reduced Linear Discriminant Analysis}

In this section we explore what happens when further reducing the dimensions and using LDA. previously we spoke about how $(K-1)$ dimension sub-space is sufficient to distinguish between $(K)$ classes, so make that $L$ now. as explained in the video, we want to maximize the spread of this $L$ dim sub-space. it means when we project this large N dim features to $L$ dim sub-space , it should retain the max information. This can be done using PCA, since that is what PCA does. \\
Note: refer to the intro to PR classes for PCA or many resources available online. 

\begin{center}
		\begin{minipage}{0.8\textwidth}
			\begin{itemize}
				\item[{\bf Problem}:] How to choose an $L$-dimensional subspace with $L\le K-1$ that is good for  LDA? 
				\item[{\bf Idea}:] Maximize the spread of the $L$-dimensional projection of centroids. 
				\item[{\bf Solution}:] Principal component analysis, i.\,e.\    we compute the principal components of the  covariance matrix of the mean vectors
				      $$\vect{\mu}'_y = \phi(\hat{\vect{\mu}}_y) \in \real^{d},$$
				      where $y=1,2,\dots, K$ and $\phi:\real^d\to\real^d$.
			\end{itemize}
		\end{minipage}
	\end{center}
		In \textbf{Principal Component Analysis (PCA)} we compute a linear mapping $\mat\Phi \in \real^{L \times d}$ that
	results in the highest spread of projected features:
	
	\begin{displaymath}
		\mat{\Phi}^* = \argmax_{\mat{\Phi}} \left(
		\frac{1}{K}
		\sum_{y=1}^K (\mat{\Phi}\vect {\mu}'_y - \mat {\Phi}\bar{\vect \mu}' )^T
		(\mat{\Phi}\vect {\mu}'_y - \mat {\Phi}\bar{\vect \mu}' ) +
		\sum_{i=1}^{L} \nu_i (\| \mat{\Phi}_i \|_2^2 - 1)
		\right)
	\end{displaymath}
	where we applied the \textbf{Lagrange multiplier} method to allow for the maximization of the spread subject to
	\begin{displaymath}
		\| \mat {\Phi}_i \|_2^2 = 1 ~, \quad i = 1, \ldots, L ~.
	\end{displaymath}	%
	Here $\| \mat{\Phi}_i \|_2^2$ denotes the squared $L_2$ norm of the $i$-th row vector of $\mat{\Phi}$, \\
	$\nu_i$ are the Lagrange multipliers, and $\bar{\vect\mu}' = \frac{1}{K}\sum_{y=1}^K\vect\mu'_y$.\\
	Note: the derivation is not straight fwd. look at intro to PR to understand the derivation. have to use fundamentals of linear algebra approach. derivation is not important for exam. 
	
	\subsection{Excursus: Optimization with Constraints}
	very brief explanation of lagrangian constraints and optimization. it is a method to find the max value that satisfies a condition relationship like below. Lagrangians is covered in more detail in optimization chapter. or look online for more resources. 
	
		\begin{itemize}
			\item Find the maximum of $f(x,y) = x +y$ 
			\item constraint: $x^2 + y^2 = 1$ 
			\item Lagrange function: \\
			      $\mathcal L(x,y,\nu) = x + y + \nu(x^2 + y^2 - 1)$ 
			\item Set the partial derivatives to zero: \\
			      $\frac{\partial \mathcal L}{\partial x} =
				      \frac{\partial \mathcal L}{\partial y} =
				      \frac{\partial \mathcal L}{\partial \nu}
				      \stackrel{!}{=} 0$
		\end{itemize}
	We need some facts from matrix calculus: 	\footnote{www.matrixcookbook.com} \footnote{website currently offline-you can still download it at \url{http://www2.imm.dtu.dk/pubdb/views/publication_details.php?id=3274} } 
	
	\begin{enumerate}
		\item Let $\vect\mu$ denote the mean and $\mat \Sigma$ the covariance matrix of a random vector $\vect x$, then we get:
		      \begin{displaymath}
			      E[(\mat A \vect x)^T(\mat A \vect x)] =
			      {\mbox{tr}}(\mat A \mat \Sigma \mat A^T) + (\mat A \vect \mu)^T(\mat A\vect \mu)
		      \end{displaymath}
		\item The matrix derivative is:
		      \begin{displaymath}
			      \frac{\partial \mbox{tr}( \mat X \mat B \mat X^T)}{\partial \mat X} = \mat X\mat B^T +\mat X\mat B
		      \end{displaymath}
	\end{enumerate}
For our \textbf{optimization problem} this implies:

	\begin{eqnarray*}
		& & \frac{\partial}{\partial{\mat{\Phi}}}
		\left\{
		\frac{1}{K} \sum_{y=1}^K (\mat{\Phi} \vect{\mu}'_y - \mat{\Phi} \bar{\vect{\mu}}' )^T
		(\mat{\Phi} \vect{\mu}'_y - \mat{\Phi} \bar{\vect{\mu}}' )  +
		\sum_{i=1}^{L} \nu_i (\| \mat{\Phi}_i \|^2_2 - 1)
		\right\} \\
		& & = \quad \frac{\partial}{\partial{\mat{\Phi}}}
		\left\{
		\frac{1}{K} \sum_{y=1}^K \left( \mat\Phi (\vect{\mu}'_y - \bar{\vect{\mu}}') \right)^T
		\left( \mat\Phi (\vect{\mu}'_y - \bar{\vect{\mu}}') \right) +
		\sum_{i=1}^{L} \nu_i (\| \mat {\Phi}_i \|^2_2 - 1)
		\right\} \\
		& & = \quad \frac{\partial}{\partial{\mat{\Phi}}}
		\left\{
		\mbox{tr} (\mat\Phi \ \mat \Sigma_{\mbox{inter}} \ \mat\Phi^T) +
		\sum_{i=1}^{L} \nu_i (\| \mat {\Phi}_i \|_2^2 - 1)
		\right\}
		\stackrel{!}{=} \mat 0,
	\end{eqnarray*}
	{\small where $\mat \Sigma_{\mbox{inter}} = \frac{1}{K}\sum_{y=1}^K (\vect \mu'_y - \bar{\vect \mu}')(\vect \mu'_y - \bar{\vect \mu}')^T$ is the inter-class (between-class) covariance matrix of the transformed class means.}
Now we compute the \textbf{partial derivatives}:

	\begin{displaymath}
		\frac{\partial}{\partial{\mat{\Phi}}}
		\left\{
		\mbox{tr} (\mat \Phi \ \mat \Sigma_{\mbox{inter}}  \ \mat \Phi^T) + \sum_{i=1}^{L} \nu_i (\| \mat {\Phi}_i \|^2_2 - 1)  \right\} =
		2 \mat \Phi\mat  \Sigma_{\mbox{inter}} +2 \mat{\Nu} \mat \Phi = \mat 0,
	\end{displaymath}
	where $\mat\Nu = \mbox{diag}(\nu_1,\dots,\nu_L)$ contains the Lagrange multipliers.
	This results in the eigenvalue and eigenvector problem:
	\begin{displaymath}
		\mat{\Sigma}_{\mbox{inter}} \mat \Phi^T = \mat \Phi^T \mat{\Lambda}', \qquad \mat\Lambda' = -\mat\Nu = \mbox{diag}(\lambda'_1,\dots,\lambda'_L),
	\end{displaymath} 
	i.\,e.\ the rows of $\mat\Phi$ are eigenvectors of $\mat\Sigma_{\mbox{inter}}$ with eigenvalues $\lambda'_i$.
	\textbf{Note:} \\
	In original PCA, the transform $\mat \Phi$ maximizes the overall spread using the covariance matrix of all features:
	\begin{displaymath}
		\mat{\Sigma}\,\mat{\Phi}^T = \mat{\Phi}^T\,\mat{\Lambda}',
	\end{displaymath}
	\begin{algorithmic}
		\STATE  \textbf{Input:} training data:  $S = \{ (\vect x_1, y_1), (\vect x_2, y_2), (\vect x_3, y_3), \dots, (\vect x_m, y_m) \}$ \\
		\STATE 1. Compute the covariance matrix of transformed mean vectors
		$$\widehat{\mat{\Sigma}}_{\mbox{inter}} =
			\frac{1}{K} \sum_{y=1}^K (\vect \mu'_y - \bar{\vect \mu}')(\vect \mu'_y - \bar{\vect \mu}')^T,$$
		\quad\,where $\bar{\vect \mu}' = \frac{1}{K} \cdot \sum_{y=1}^{K}\vect\mu'_y$. \\
		\STATE 2. Compute the $L$ eigenvectors of the covariance matrix belonging \\
		\quad\,to the largest eigenvalues. \\
		\STATE 3. The eigenvectors are the rows of the mapping $\mat\Phi\in\real^{L\times d}$.
		\quad\,(The class centroids span at most a $(K-1)$-dimensional subspace of $\real^d$.) \\
		\STATE  \textbf{Output:} matrix $\mat\Phi$
	\end{algorithmic}
	
	\subsection{Fisher Transform}%
	References to read are given below \footnote{\url{https://www.stat.cmu.edu/~ryantibs/datamining/lectures/21-clas2.pdf}}%
	\footnote{\url{https://www.cs.cmu.edu/~psarkar/sds383c_16/lecture_scribe13.pdf}}.

	The described method to compute the LDA mapping is not the original derivation. Some additional reading material.\autocite{fisher1936,bishop2006}%
	\footnote{\url{https://towardsdatascience.com/fishers-linear-discriminant-intuitively-explained-52a1ba79e1bb}}%
	\footnote{\url{https://www.ccis.northeastern.edu/home/vip/teach/MLcourse/5_features_dimensions/lecture_notes/LDA/LDA.pdf}}

\textbf{Original method}
		\begin{itemize}
			\item Project samples $\vect{x}_i$ onto a straight line with direction $\vect{r}$, $\lVert \vect{r} \rVert_2 = 1$:
			      \begin{displaymath}
				      x_i^{\mathrm{proj}} = \vect{x}_i^T \vect{r}
			      \end{displaymath}
			\item Maximize the ratio of the between-class scatter and the within-class scatter:
			      \begin{displaymath}
				      \vect{r}^* =
				      \argmax_{\vect{r}} R(\vect{r}) =
				      \argmax_{\vect{r}} \frac{|\mu_1^{\mathrm{proj}} - \mu_2^{\mathrm{proj}}|^2}{(s_1^{\mathrm{proj}})^2 + (s_2^{\mathrm{proj}})^2}
			      \end{displaymath}
			\item Classify by applying a threshold to $x_i^{\mathrm{proj}}$
		\end{itemize}
	\textbf{Finding $\vect{r}^*$}

	\begin{enumerate}
		\item Mean and scatter matrix for each class:
		      {\small
		      \begin{eqnarray*}
			      \vect{\mu}_k &=& \frac{1}{m_k} \sum_{i:\, y_i = k} \vect{x}_i \\
			      \mat{S}_k   &=& \sum_{i:\, y_i = k} (\vect{x}_i - \vect{\mu}_k) (\vect{x}_i - \vect{\mu}_k)^T
		      \end{eqnarray*}
		      }
		\item \textbf{Within-class scatter} matrix:
		      \begin{displaymath}
			      \mat{S}_{\mathsf{W}} = \mat{S}_1 + \mat{S}_2
		      \end{displaymath}
		\item \textbf{Between-class scatter} matrix:
		      \begin{displaymath}
			      \mat{S}_{\mathsf{B}} = (\vect{\mu}_1 - \vect{\mu}_2) (\vect{\mu}_1 - \vect{\mu}_2)^T
		      \end{displaymath}
	
		\item Expressing $\mu_k^{\mathrm{proj}}$ and $(s_k^{\mathrm{proj}})^2$ of the projected samples in terms of $\vect{\mu}_k$ and $\mat{S}_k$:
		      {\small
		      \begin{eqnarray*}
			          \mu_k^{\mathrm{proj}} &=&
				         \frac{1}{m_k} \sum_{i:\, y_i = k} x_i^{\mathrm{proj}} 
				       = \frac{1}{m_k} \sum_{i:\, y_i = k} \vect{r}^T  \vect{x}_i 
				       = \vect{r}^T \vect{\mu}_k  \\
				      (s_k^{\mathrm{proj}})^2 &=&
				        \sum_{i:\, y_i = k} (x_i^{\mathrm{proj}} - \mu_k^{\mathrm{proj}})^2 
				      = \sum_{i:\, y_i = k} (\vect{r}^T \vect{x}_i - \vect{r}^T \vect{\mu}_k)^2 
				      = \vect{r}^T \mat{S}_k \vect{r} 
		      \end{eqnarray*}
		      }
		\item Plug it into $R(\vect{r})$:
		      {\small
		      \begin{displaymath}
			      R(\vect{r})
			      = \frac{|\mu_1^{\mathrm{proj}} - \mu_2^{\mathrm{proj}}|^2}{(s_1^{\mathrm{proj}})^2 + (s_2^{\mathrm{proj}})^2}
			      = \frac{\vect{r}^T \mat{S}_B \vect{r}}{\vect{r}^T \mat{S}_W \vect{r}}
			      \end{displaymath}
		      }
		\item[] This is known as the \textbf{Generalized Rayleigh Quotient}.
		\item \textbf{Why this yields an eigenproblem:} since $R(\vect r)$ does not change if $\vect r$ is scaled, we can fix the denominator and solve
		      {\small
		      \[
		        \max_{\vect r}\; \vect r^T\mat S_\mathsf B\vect r
		        \quad\text{subject to}\quad \vect r^T\mat S_\mathsf W\vect r = 1~.
		      \]
		      }
		      With Lagrangian $\Lambda(\vect r,\lambda)=\vect r^T\mat S_\mathsf B\vect r-\lambda(\vect r^T\mat S_\mathsf W\vect r-1)$, stationarity gives
		      $\mat S_\mathsf B\vect r = \lambda\,\mat S_\mathsf W\vect r$.
		      Moreover, multiplying from the left by $\vect r^T$ yields $R(\vect r)=\lambda$, hence we take the generalized eigenvector with the \textbf{largest} eigenvalue.
		\item Maximizing the Generalized Rayleigh Quotient is equivalent to solving the following \textbf{generalized eigenvalue problem}:
		      {\small
		      \begin{eqnarray*}
			      \mat{S}_\mathsf{B} \vect{r}^* &=& \lambda \mat{S}_\mathsf{W} \vect{r}^* \\
			      \mat{S}_\mathsf{W}^{-1} \mat{S}_\mathsf{B} \vect{r}^* &=& \lambda \vect{r}^* \\
		      \end{eqnarray*}
		      }
		\item \textbf{Note}: $\mat{S}_\mathsf{B} \vect{r}^*$ is always in the direction of $\vect{\mu}_1 - \vect{\mu}_2$; \\
		      no need to compute the eigenvalues and eigenvectors of $\mat{S}_\mathsf{W}^{-1} \mat{S}_\mathsf{B}$! \\
		      The direction of $\vect{r}^*$ is:
		      {\small
		      \begin{displaymath}
			      \vect{r}^* = \mat{S}_\mathsf{W}^{-1} (\vect{\mu}_1 - \vect{\mu}_2)
		      \end{displaymath}
		      }	      
	\end{enumerate}
	
	\begin{itemize}
		\item Usually the total linear mapping for LDA is computed dimension by dimension through the maximization of the Rayleigh ratio for each projection axis $\vect a^*$:
		      \begin{eqnarray*}
			      \vect a^* &=& \argmax_{\vect a}\ \frac{\vect a^T \mat \Sigma_{\mbox{inter}} \vect a}{\vect a^T \mat \Sigma_{\mbox{intra}} \vect a}
		      \end{eqnarray*}
		\item The solution is a generalized eigenvalue problem: $\vect a$ is the eigenvector of
		      $$\mat \Sigma_{\mbox{intra}} ^{-1} \mat \Sigma_{\mbox{inter}}$$
		      that belongs to the largest eigenvalue.
	\end{itemize}
	In literature the optimization problem is mostly rewritten:
		\begin{itemize}
		\item Equivalent constrained optimization problem
		      \begin{eqnarray*}
			      \mbox{maximize:}   & & \vect r^T \mat \Sigma_{\mbox{inter}}\vect r\\
			      \mbox{subject to:} & & \vect r^T \mat \Sigma_{\mbox{intra}} \vect r =1
		      \end{eqnarray*}
		\item Lagrange multiplier method ($\lambda>0$):
		      \begin{displaymath}
			      \vect{r}^* = \argmax_{\vect r}
			      \left\{
			      \vect r^T \mat \Sigma_{\mbox{inter}} \vect r
			      - \lambda\big(\vect r^T \mat \Sigma_{\mbox{intra}} \vect r - 1\big)
			      \right\}
		      \end{displaymath}
	\end{itemize}

In a nutshell, we recognize the fishers problem is solved by finding the eigenvalues. 

\subsection{Dimensionality Reduction}


	A few comments on dimensionality reduction:

	\begin{itemize}
		\item PCA does not require a classified set of feature vectors  (in contrast to LDA). 
		\item PCA transformed features are approximately normally distributed (central limit theorem).
		\item Components of PCA transformed features are mutually uncorrelated. *Independent and uncorrelated if gaussian distribution. 
		\item There exist many other methods for dimensionality reduction, e.\,g.,  Sammon transform, independent component analysis. 
		\item Usually the estimation of transforms is computationally prohibited
		\item \textbf{Johnson-Lindenstrauss lemma}: If vectors are projected onto a randomly selected subspace of suitably high dimension, then the distances between the vectors are approximately preserved.
	\end{itemize}
	
\subsection{LDA application: adidas\_1 }
Refer to the slides for more detailed information. The key takeaway is to understand how a practical classification problem was solved and what was the output of using an LDA classifier. 

\subsection{PCA application: Shape Modeling}

Details in the slides, but the key takeaway is to understand how PCA is used for shape modeling. it can be easily understood that the mean of all the features represents the average shape of the organ, and PCA helps us to understand the amount of variation seen in the data. 

\subsection{Notes on Regression}
We remember that regression is basically trying to estimate a sample, not just understanding which class it belongs to. In the example of apples vs pears, we are trying to estimate the shape, size, weight of the apple using the knowledge of probabilities, prior class information about how apples are usually shaped, weigh, etc.\autocite{hastie2009,bishop2006}

\subsubsection{Linear Regression}
In the two class situation, we set $y\in \{-1,+1\}$ and use the decision rule:
	\begin{displaymath}
		\hat y = \mbox{sgn}(\vect \alpha^T\vect x + \alpha_0).
	\end{displaymath}
	We can compute the linear decision boundary simply by least-square estimation. For a given set of learning data we use matrix notation:

	\begin{displaymath}
		\mat{X} = \left(
		\begin{array}{cc}
				\vect x_1^T & 1      \\
				\vect x_2^T & 1      \\
				\vdots     & \vdots \\
				\vect x_m^T & 1
			\end{array}
		\right) \in \real^{m\times (d+1)} \qquad \vect y \quad = \quad
		\left(
		\begin{array}{c}
				y_1    \\
				y_2    \\
				\vdots \\
				y_m    \\
			\end{array}
		\right) \in \real^m
	\end{displaymath}
	and define
	\begin{displaymath}
		\vect \theta = { \vect \alpha\choose \alpha_0} \in \real^{d+1}, \qquad
		\tilde{\vect x}_i = { \vect x_i \choose 1} \in \real^{d+1}.
	\end{displaymath}
	i.\,e.\ the $i$-th row of $\mat X$ is the augmented feature vector $\tilde{\vect x}_i^T$ and $\vect\theta^T\tilde{\vect x}_i = \vect\alpha^T\vect x_i + \alpha_0$.

		One option to estimate $\vect \theta$ is to solve the linear regression problem:
	\begin{displaymath}
		\hat{\vect \theta} = \argmin_{\vect \theta}\| \mat X \vect \theta - \vect y\|_{2}^2
	\end{displaymath}
The \textbf{}{least-square estimator} for the $L_2$-norm:
	\begin{eqnarray*}
		\widehat{\vect \theta} &=& \argmin_{\vect\theta} \ \sum_{i=1}^m (\vect \theta^T\tilde{\vect x}_i-y_i)^2\\
		&=& \argmin_{\vect\theta} \ (\mat{X}\vect \theta-\vect y)^T(\mat{X}\vect \theta-\vect y)
	\end{eqnarray*}
	and thus we get
	\begin{displaymath}
		\widehat{\vect \theta} = (\mat{X}^T\mat{X})^{-1}\mat{X}^T\vect y
	\end{displaymath}
	if the column vectors of $\mat X$ are linearly independent. A few obvious questions:
	\begin{itemize}
		\item Why should we prefer the Euclidean norm ($L_2$-norm)? 
		\item Will different norms lead to different results? 
		\item Which norm and decision boundary is the best one? 
		\item Can we incorporate prior knowledge in linear regression?
	\end{itemize}	

\subsubsection{Ridge Regression}

A regression formula based on L2 norm and the regularizer is also an L2 norm. regularized expression helps to make sure, that we are always able to find the inverse. since the solution depends on finding the inverse of a matrix. \\ 

\textbf{Regularized regression:}

	\begin{eqnarray*}
		\widehat{\vect\theta}
		&=& \argmin_{\vect\theta} \| \mat{X} \vect\theta - \vect y \|^2_2 + \lambda \|\vect\theta\|^2_2 \\ 
		&=& \argmin_{\vect\theta} \ (\mat{X} \vect\theta - \vect y)^T (\mat{X} \vect\theta-\vect y) +
		\lambda \cdot \vect\theta^T\vect \theta
	\end{eqnarray*}
	and thus we get the estimator:
	\begin{displaymath}
		\widehat{\vect\theta} = (\mat{X}^T\mat{X}+\lambda \mat{I})^{-1}\mat{X}^T\vect y
	\end{displaymath}
	\textbf{Notes:}
	\begin{itemize}
		\item The term $\lambda \mat{I}$ adds a positive constant $\lambda$ to the diagonal elements.
		\item The problem is non-singular even if $\mat{X}^T \mat{X}$ is not of full rank. 
		\item This was the main motivation of ridge regression when it was first introduced in statistics in 1970.
		\item The ridge solutions are not equivariant under scaling of the inputs: 
		      \textbf{standardize} the input before solving the regression problem!  
		\item The intercept $\alpha_0$ should not be penalized, there should be bias correction which is done as shown below:
		      \begin{itemize}
			      \item \textbf{Center} the input $\vect{x}_i$.
			      \item Estimate $\alpha_0$ by $\bar{y} = \frac{1}{m} \sum_{i=1}^m y_i$.
			      \item Estimate the remaining coefficients by a ridge regression without  intercept. Matrix $\vect{X}$ has $d$ columns (instead of $d+1$).
		      \end{itemize}
	\end{itemize}
\textbf{Statistical approach:} parameters $\alpha_j$ are random variables 
	\begin{itemize}
		\item It is extended by a prior distribution on the parameter vector $\vect\theta$:
		      \begin{displaymath}
			      \vect \theta \sim {\cal N}(\vect 0, \tau^2 \mat I)
		      \end{displaymath}
		      (in practice the intercept $\alpha_0$ is excluded from the prior/penalty, see below)
		\item Suppose
		      \begin{displaymath}
			      \forall 1 \le i \le m: \quad y_i
			      \sim \mathcal{N}(\underbrace{\vect{\alpha}^T \vect{x}_i + \alpha_0}_\text{mean}, \underbrace{\sigma^2}_\text{variance})
			      ,\ \text{i.e. } p(y_i|\vect x_i,\vect\alpha) = \frac{1}{\sqrt{2 \pi} \cdot \sigma} e^{-\frac{1}{2} \cdot \frac{(y_i - \vect{\alpha}^T \vect{x}_i - \alpha_0)^2}{\sigma^2}}
  \end{displaymath}
		\item Parameters $\alpha_j$ are assumed to be independent of each other.
		\item Prior distribution of $\alpha_j$:
		      \begin{displaymath}
			      \forall 1 \le j \le d: \quad \alpha_j
			      \sim \mathcal{N}(0, \tau^2)
			      ,\ \text{i.e. } p(\alpha_j) = \frac{1}{\sqrt{2 \pi} \cdot \tau} e^{-\frac{1}{2} \cdot \frac{(\alpha_j - 0)^2}{\tau^2}}
		      \end{displaymath}
		\item {Maximizing the posterior probability of $\vect{\alpha}$ for given $\sigma^2$ and $\tau^2$:}
		      \footnotesize
		      \begin{eqnarray*}
			      \argmax_{\vect{\alpha}} p(\vect{\alpha}|S) 
			      &=& \argmax_{\vect{\alpha}} \Bigg\{ p(\vect{\alpha}) \cdot \prod_{i=1}^m p(y_i | \vect x_i, \vect{\alpha}) \Bigg\} \\ 
			      &=& \argmax_{\vect{\alpha}} \Bigg\{ \prod_{j=1}^d p(\alpha_j) \cdot \prod_{i=1}^m p(y_i | \vect x_i, \vect{\alpha}) \Bigg\} \\ 
			      &=& \argmax_{\vect{\alpha}} \Bigg\{ \sum_{j=1}^d \log p(\alpha_j) + \sum_{i=1}^m \log p(y_i | \vect x_i, \vect{\alpha}) \Bigg\} \\ 
			      &=& \argmax_{\vect{\alpha}} \Bigg\{ - \frac{1}{2 \tau^2} \sum_{j=1}^d \alpha_j^2 - \frac{1}{2 \sigma^2} \sum_{i=1}^m (y_i - \vect{\alpha}^T \vect{x}_i - \alpha_0)^2 \Bigg\} \\ 
			      &=& \argmin_{\vect{\alpha}} \Bigg\{ \frac{\sigma^2}{\tau^2} \sum_{j=1}^d \alpha_j^2 + \sum_{i=1}^m (y_i - \vect{\alpha}^T \vect{x}_i - \alpha_0)^2 \Bigg\}
		      \end{eqnarray*}
    		\item {Maximizing the posterior probability of $\vect{\alpha}$ for given $\sigma^2$ and $\tau^2$:} 
		      {\footnotesize
		      \begin{displaymath}
			      \argmax_{\vect{\alpha}} p(\vect{\alpha}|S)
			      = \argmin_{\vect{\alpha}} \Bigg\{ \lambda \vect{\alpha}^T \vect{\alpha} + (\mat{X}\vect{\theta} - \vect{y})^T (\mat{X}\vect{\theta} - \vect{y}) \Bigg\}
			      \quad \text{with} \quad
			      \lambda = \frac{\sigma^2}{\tau^2}
		      \end{displaymath}
		      where $\vect\theta = (\vect\alpha^T, \alpha_0)^T$ (the intercept $\alpha_0$ is not penalized).
		      }
		\item {The ridge estimate is the mode of the posterior pdf!}
	\end{itemize}

\subsubsection{Lasso}

A different kind of regression that is based on L1 norm regularizer. 

	Regularized regression using a mixture of $L_2$- and $L_1$-norm, where the residual is penalized using the $L_2$-norm and the regularizer uses the $L_1$-norm:

	\begin{displaymath}
		\widehat{\vect \theta} = \argmin_{\vect\theta} \|\mat{X}\vect \theta-\vect y\|^2_2 + \lambda \cdot \|\vect \theta\|_1
	\end{displaymath}
	The lasso is used to compute a sparse solution of the system of linear equations, i.\,e.\ the number of non-zero elements in $\vect \theta$ shall be small.
	
\subsection{Lessons Learned}

	\begin{itemize}
		\item Principal component analysis 
		\item Linear discriminant analysis with and without dimension reduction 
		\item Both PCA and LDA relate to an eigenvalue eigenvector problem 
		\item Alternative formulation of LDA using the Fisher transform 
		\item Linear and ridge regression for classification
	\end{itemize}
	
\subsection{Comprehensive Questions}
	\begin{itemize}
		\item What is the difference between PCA and LDA? 
		\item How can PCA and LDA be combined to achieve a high rank reduction? 
		\item Write down a straight forward objective function for linear regression! 
		\item What happens if we replace the $L_2$-norm by another norm?
	\end{itemize}
	
